\section{Prior methods, and how each one fails here}
\label{sec:related}

\begin{lede}
Four families of solution, in increasing order of how much they give up for
speed. Each is described by what it optimises, what one timestep costs it, and
the specific way it breaks in a narrow lifelong warehouse --- because that
failure mode is what the next section is designed around.
\end{lede}

\subsection{Optimal MAPF: CBS and its relatives}
\label{sec:related-cbs}

Conflict-Based Search \citep{sharon2015cbs} solves one-shot MAPF optimally with
a two-level search: a low level plans each agent independently, and a high
level branches on the first conflict found between two agents' paths, adding a
constraint forbidding one of them from the conflicting vertex or edge at that
time and replanning only that agent. It is complete and optimal for sum-of-costs,
and it is exponential in the number of conflicts --- which in a dense warehouse
corridor is very nearly the number of agent pairs.

Two properties disqualify it as-is here. First, cost: the high-level tree grows
with contention, and contention is exactly the condition of interest. Second,
and more fundamental, CBS answers a one-shot question. In a lifelong setting
goals change continuously, so the plan CBS produces is stale before it is
finished. Every practical lifelong system built on CBS therefore truncates it,
which is what RHCR does.

\subsection{PIBT: one timestep, decided locally}
\label{sec:related-pibt}

Priority Inheritance with Backtracking \citep{okumura2019pibt} abandons
multi-step planning entirely. At every timestep, robots are considered in
descending priority order; each takes the best-scoring legal move available to
it. If the cell it wants is occupied by a robot that has not yet decided, the
requesting robot \emph{lends its priority} to the occupant and asks it to move
first --- recursively. If the occupant can find no legal move, the request
backtracks and the requester tries its next candidate. \Cref{sec:pibt} gives
the recursion in full.

Two properties matter for everything that follows.

\begin{enumerate}
\item \textbf{Cost.} One timestep is $O(|A|(\Delta(G) + F + \log|A|))$ for a
fleet $A$, maximum degree $\Delta(G)$ and per-candidate scoring cost $F$. On
the maps used here this is 0.4--1.0\,ms per step, one to two orders of
magnitude below either baseline in \Cref{sec:results}.

\item \textbf{The push property.} Priority inheritance works because a robot
asked to move can, in general, be pushed somewhere. PIBT's progress guarantee
is proved under the assumption that every edge of $G$ lies on a simple cycle;
informally, a robot displaced from a cell always has an alternative. This
assumption is not a technicality --- \Cref{sec:contribution-soft} shows that
the natural way to implement one-way aisles \emph{violates it by construction},
and that this, rather than anything about direction control as an idea, was the
cause of every collapse previously attributed to the aisle layer.
\end{enumerate}

What PIBT does not have is any notion of structure above a single cell. It
cannot know that a corridor is a corridor, cannot commit a corridor to one
direction, and cannot coordinate two robots that will meet in ten timesteps.
Every robot's decision is local and one step deep. That is the gap this work
addresses.

\subsection{Token Passing: reserve a path, or wait}
\label{sec:related-tp}

Token Passing \citep{ma2017tp} is the standard lifelong-MAPD baseline for
exactly this pickup-and-delivery setting. A token circulates; a robot holding
it plans a collision-free path to its goal through a shared reservation table
using cooperative A$^\ast$ over $(\text{vertex}, \text{time})$ states, commits
that path into the table, and passes the token on. A robot that cannot find a
path waits in place.

Its failure mode in a narrow warehouse is structural, not a matter of tuning:
\textbf{Token Passing has no mechanism by which a blocked robot can displace
the robot blocking it}. Once robots queue nose-to-tail in a single-file
corridor, no robot in the queue can reserve a path through the robots ahead of
it, every robot's search fails, every robot waits, and the configuration is
absorbing --- it never resolves for the rest of the run. On
\variant{warehouse\_bottleneck}, whose two halves are joined by one corridor,
all 16 robots reach this state and the run completes essentially no tasks. The
first comparison animation in \code{docs/gifs/} shows precisely this, beside
the same scenario under priority inheritance.

This is the failure PIBT was invented to remove, and it is worth being explicit
that observing it here is a confirmation of the published analysis rather than
a criticism of Token Passing: the algorithm was never claimed to be live in
maps without alternative routes.

\subsection{RHCR: rolling-horizon windowed CBS}
\label{sec:related-rhcr}

Rolling-Horizon Collision Resolution \citep{li2021rhcr} makes CBS lifelong by
bounding it in time. Every \param{replan\_period} timesteps all robots are
jointly replanned over a window of \param{window} timesteps using CBS
restricted to that window; conflicts beyond the window's edge are simply not
resolved yet, on the reasoning that they will be replanned before they arrive.
Between replans, individual stale paths are repaired.

RHCR is the strongest of the three published methods for this setting and its
design target is large: hundreds of robots on warehouses far bigger than the
ones used here, where a rolling horizon pays for itself against full CBS. At
the scale measured in this document, with untuned window parameters
($10/5$), it costs 43--98\,ms per timestep --- against 0.4--4.4\,ms for the
PIBT variants --- and does not convert that cost into throughput. We report
this as a property of \emph{this baseline at this scale with these defaults}
and not as a claim about RHCR generally; \Cref{sec:results-caveats} states the
caveat again where the numbers appear.

\subsection{Directional and structural traffic control}
\label{sec:related-directional}

Making corridors one-way is not a new idea; warehouse floors have been run that
way by human traffic rules for as long as there have been warehouses, and
highway-style directional layouts appear throughout the MAPF and AGV
literature. The design decisions that distinguish proposals are: at what
granularity a direction is decided (a whole map, a corridor, an edge), how
often it may change, and --- the question this work turns out to hinge on ---
\textbf{what a committed direction does to a move that opposes it}.

The usual answer is that it forbids it. A one-way corridor is one-way; a move
against it is not legal; the planner never sees it. Under a planner that plans
whole paths, this is harmless: the path simply routes around. Under PIBT it is
not harmless at all, for the reason given in \Cref{sec:related-pibt}, and that
is the subject of the next section.
